% TEMPLATE-MILC-v3.tex - plantilla maestra UNICA de las guias MILC v3.
%
% La usan las cuatro familias de guias, cada una con su driver:
%   · Grados 6-11          -> scripts/build-guias-g11.py
%   · Bebras               -> scripts/build-guias-bebras.py
%   · Semillero            -> scripts/build-guias-semillero.py
%   · Territorio interior  -> scripts/build-guias-territorio-interior.py
%
% Antes habia tres copias casi identicas (~400 lineas iguales, 6 distintas):
% cada mejora habia que aplicarla tres veces, y en la practica no se hacia
% (el motor de recursos vivia solo en dos de ellas). Lo que varia entre
% familias esta parametrizado en 6 placeholders que cada driver llena
% (se nombran aqui SIN su sintaxis real: el builder reemplaza por texto plano,
% tambien dentro de los comentarios, y un valor multilinea rompe el preambulo):
%
%   HEADER_IZQ        encabezado izquierdo de pagina
%   HEADER_DER        encabezado derecho de pagina
%   PORTADA_ETIQUETA  rotulo sobre el numero grande (GRADO/MOMENTO/MODULO)
%   PORTADA_SERIE     linea de serie de la portada
%   PORTADA_ID        identificador de la guia en la portada
%   PORTADA_CIRCULO   el circulito de grado (°), o vacio si no aplica
%   PDF_TITULO        titulo en texto plano (sin \textbf ni comillas TeX) para
%                     los metadatos del PDF (/Title); lo produce lib_xelatex.texto_pdf
%
% Composicion (auditoria 2026-09-03): las bandas de estacion piden espacio
% (\Needspace*) y prohiben el salto justo despues (\nopagebreak), asi no quedan
% huerfanas al pie; las cajas largas (softbox, drawbox, infoband, puentebox) son
% `breakable`, asi una caja de media pagina ya no arrastra todo a la siguiente
% dejando la anterior al 40 %. Todo texto sobre mostaza o verde va en milcNegro
% (blanco sobre #E5B400 da 1,93:1 y sobre #58A923 2,95:1; ninguno pasa WCAG AA).
% El resto de claves las documenta content/guias/_SCHEMA.md. El script valida
% que no queden placeholders sin reemplazar antes de escribir el archivo final.

% Etiquetado PDF (latex-lab): /Lang, estructura y texto alternativo de las
% figuras, para lectores de pantalla. Debe ir ANTES de \documentclass.
\DocumentMetadata{lang=es-CO,tagging=on}
\documentclass[11pt,letterpaper]{article}
% headheight=24pt: la cabecera derecha lleva dos lineas (identidad de la guia
% y estacion actual). headsep se acorta en la misma medida para que el bloque
% de texto quede exactamente donde estaba con headheight=12pt.
\usepackage[margin=2.0cm,headheight=24pt,headsep=13pt]{geometry}
\usepackage[table]{xcolor}
\usepackage{fontspec}
\usepackage[spanish]{babel}
\usepackage{tikz}
% Librerías TikZ para los diagramas de las guías (bloque `recursos:`):
%  · babel  → babel en español vuelve activos caracteres como «>», y TikZ los usa
%             en sus opciones (>=stealth, ->). Sin esta librería, cualquier
%             diagrama con flechas revienta con «\language@active@arg> has an
%             extra }». Debe ir ANTES de que se dibuje nada.
%  · positioning        → sintaxis «below left=2pt and 3pt».
%  · decorations.pathreplacing → llaves ({brace}) para señalar rangos.
\usetikzlibrary{babel,positioning,decorations.pathreplacing}
\usepackage{graphicx}
% Circuitos: el trabajo de electrónica se hace hoy con micro:bit y sus sensores
% integrados (MakeCode, sin cableado), así que no se necesitan esquemáticos.
% Si algún día entran montajes con protoboard, basta descomentar esta línea:
% los diagramas .tex/.tikz declarados en `recursos:` ya se incrustan solos.
% \usepackage{circuitikz}
\usepackage[most]{tcolorbox}
\usepackage{tabularx}
\usepackage{array}
\usepackage{fancyhdr}
\usepackage{titlesec}
\usepackage[document]{ragged2e}  % bandera a la derecha en todo el cuerpo
\usepackage{enumitem}
\usepackage{microtype}
\usepackage{etoolbox}
\usepackage{needspace}
% hyperref al final del preambulo: metadatos del PDF (titulo, autor, idioma,
% familia), marcadores por estacion (\pdfbookmark) y enlaces sin recuadro.
\usepackage[hidelinks,
  pdftitle={Sustentación del proyecto — decir qué funciona y qué no},
  pdfauthor={PhD. Álvaro Cárdenas Orozco},
  pdfsubject={Tecnología e Informática · Grado 8 · Guía 20 · MILC v3},
  pdflang={es-CO},
  bookmarksopen=true,bookmarksnumbered=false]{hyperref}

% Helvetica Neue no trae la flecha U+2192, asi que cada «→» escrito literalmente
% en el YAML se compilaba a NADA, en silencio (462 flechas en 61 guias: rutas del
% tipo «paleta → tipografia → lockup» salian rotas). Mapeamos el caracter a su
% version matematica, que si tiene glifo. Asi el autor puede seguir escribiendo
% «→» comodamente en el YAML.
\usepackage{newunicodechar}
\newunicodechar{→}{\ensuremath{\rightarrow}}
% Mismo problema con los simbolos de marca y vinetas: el «✓» de «una palomita ✓»
% salia como cuadrito vacio en el PDF de todas las guias que lo usaban.
\usepackage{pifont}
\newunicodechar{✓}{\ding{51}}
\newunicodechar{✗}{\ding{55}}
\newunicodechar{✔}{\ding{52}}
\newunicodechar{•}{\textbullet}
\newunicodechar{≥}{\ensuremath{\geq}}
\newunicodechar{≤}{\ensuremath{\leq}}

\setmainfont{Helvetica Neue}
\setsansfont{Helvetica Neue}
\setlength{\parindent}{0pt}
\setlength{\parskip}{6pt}
% Sin particion de palabras: en 6.º-7.º una palabra cortada por guion al final
% de linea («Comuni-cacion») frena la lectura mucho mas que una linea corta.
\hyphenpenalty=10000
\exhyphenpenalty=10000
\pretolerance=2000
\tolerance=3000
\setlength{\emergencystretch}{3em}
\linespread{1.10}
\renewcommand{\arraystretch}{1.25}
\setlist{leftmargin=5mm,itemsep=1mm,topsep=1mm}
\newcolumntype{Y}{>{\RaggedRight\arraybackslash}X}

\definecolor{milcMagenta}{HTML}{E6007E}
\definecolor{milcVino}{HTML}{5A0038}
\definecolor{milcTurquesa}{HTML}{0093A5}
\definecolor{milcVerde}{HTML}{58A923}
\definecolor{milcMostaza}{HTML}{E5B400}
\definecolor{milcOcre}{HTML}{B86600}
\definecolor{milcNegro}{HTML}{241816}
\definecolor{milcGris}{HTML}{F7F7F4}
\definecolor{milcLinea}{HTML}{E8E2DD}
\definecolor{milcGradePrimary}{HTML}{FF6600}
\definecolor{milcGradeDark}{HTML}{9A3E00}
\definecolor{milcGradeSoft}{HTML}{FFE4D0}
\definecolor{milcGradeAccent}{HTML}{CFE7FF}
\definecolor{milcGradeText}{HTML}{FFFFFF}
\color{milcNegro}

\pagestyle{fancy}
\fancyhf{}
% Cabecera de dos lineas: identidad de la guia arriba y, a la derecha y en
% gris, la estacion en curso (\markright desde \milcestacion). Antes de la
% primera estacion la segunda linea va vacia; el \strut mantiene la altura
% para que la linea de arriba no baile entre paginas.
\lhead{\small Tecnología e Informática · Grado 8\\[-1pt]{\footnotesize\strut}}
\rhead{\small Guía 20 · MILC v3\\[-1pt]{\footnotesize\strut\color{milcNegro!65}\rightmark}}
\cfoot{\small\thepage}
\renewcommand{\headrulewidth}{0pt}

% Sobre mostaza (#E5B400) y verde (#58A923) el texto va en milcNegro; sobre el
% resto de la paleta, en blanco. Se usa dentro de fonttitle/fontupper, donde un
% \color posterior manda sobre coltitle.
\newcommand{\milctintasobre}[1]{%
  \ifboolexpr{test{\ifstrequal{#1}{milcMostaza}} or test{\ifstrequal{#1}{milcVerde}}}%
    {\color{milcNegro}}{\color{white}}}

\newtcolorbox{titlebox}[1]{
  enhanced,colback=#1,colframe=#1,arc=7mm,boxrule=0pt,
  left=7mm,right=7mm,top=3mm,bottom=3mm,
  before skip=8pt,after skip=8pt,
  fontupper=\sffamily\bfseries\Large\color{white}
}

\newtcolorbox{softbox}[3]{
  enhanced,breakable,pad at break=3mm,
  colback=#2,colframe=#1,borderline west={4pt}{0pt}{#1},
  arc=2mm,boxrule=.4pt,left=4mm,right=4mm,top=3mm,bottom=3mm,
  before skip=6pt,after skip=8pt,title={#3},coltitle=white,
  colbacktitle=#1,fonttitle=\sffamily\bfseries\milctintasobre{#1},fontupper=\RaggedRight,
  attach boxed title to top left={xshift=3mm,yshift=-2mm},
  boxed title style={arc=2mm, boxrule=0pt}
}

\newtcolorbox{drawbox}[2]{
  enhanced,breakable,pad at break=4mm,
  colback=white,colframe=#1,arc=4mm,boxrule=1pt,
  left=5mm,right=5mm,top=4mm,bottom=4mm,before skip=8pt,after skip=8pt,
  title={#2},coltitle=white,colbacktitle=#1,fonttitle=\sffamily\bfseries\milctintasobre{#1},
  fontupper=\RaggedRight,
  attach boxed title to top left={xshift=4mm,yshift=-2mm},
  boxed title style={arc=3mm, boxrule=0pt}
}

\newtcolorbox{infoband}[1]{
  enhanced,breakable,pad at break=2mm,colback=#1!8,colframe=#1!50,arc=2mm,boxrule=.5pt,
  left=4mm,right=4mm,top=2mm,bottom=2mm,before skip=4pt,after skip=4pt,
  fontupper=\small\RaggedRight
}

% Puente narrativo entre secciones (contrato editorial Conéctate).
% Da continuidad: "Acabas de X. Ahora vas a Y porque Z."
% Visualmente: banda izquierda en color del grado, texto en itálica pequeña.
\newtcolorbox{puentebox}{
  enhanced,breakable,pad at break=2mm,colback=milcGradeSoft,colframe=milcGradeDark,
  borderline west={3pt}{0pt}{milcGradeDark},
  arc=0mm,boxrule=0pt,left=5mm,right=4mm,top=2.5mm,bottom=2.5mm,
  before skip=4pt,after skip=8pt,
  fontupper=\itshape\small\color{milcNegro}\RaggedRight
}
% La flecha va en modo matematico a proposito: Helvetica Neue NO tiene el glifo
% U+2192, asi que \textrightarrow se compilaba a NADA («Missing character: There
% is no → in font Helvetica Neue») y el puente salia sin flecha. En modo
% matematico la flecha viene de la fuente matematica y si se dibuja.
\newcommand{\puente}[1]{%
  \begin{puentebox}%
    \textcolor{milcGradeDark}{$\rightarrow$}\hspace{2mm}#1%
  \end{puentebox}%
}

\newcommand{\linead}[1]{\noindent\color{milcLinea}\rule{#1}{0.5pt}\color{milcNegro}}
\newcommand{\respuesta}[1]{\par\vspace{2mm}\foreach \n in {1,...,#1}{\noindent\color{milcLinea}\rule{\linewidth}{0.5pt}\par\vspace{5mm}}\color{milcNegro}}
\newcommand{\checkbox}{\textcolor{milcGradeDark}{\fbox{\phantom{x}}}\hspace{5pt}}

% ─── Listas aireadas para las guías socioemocionales ────────────────────
% Componer las verificaciones como «\checkbox texto\\» deja el texto que salta
% de línea debajo del cuadrito y apelmaza el bloque. Estos entornos alinean la
% continuación y airean los ítems. Los usa Territorio Interior; cualquier
% familia puede hacerlo.
\newlist{checklista}{itemize}{1}
\setlist[checklista]{
  label=\textcolor{milcGradeDark}{\fbox{\phantom{x}}},
  leftmargin=9mm, labelsep=3.2mm, itemsep=2.4mm,
  topsep=2.5mm, parsep=0pt, partopsep=0pt, before=\RaggedRight
}
\newlist{pasoslista}{enumerate}{1}
\setlist[pasoslista]{
  label=\textcolor{milcGradeDark}{\sffamily\bfseries\arabic*.},
  leftmargin=8mm, labelsep=3mm, itemsep=2.8mm,
  topsep=1mm, parsep=0pt, partopsep=0pt, before=\RaggedRight
}

\newcommand{\phasecard}[4]{
\begin{tcolorbox}[enhanced,colback=#3,colframe=#2,arc=3mm,boxrule=.5pt,height=3.05cm,valign=center,halign=center,left=1.5mm,right=1.5mm,top=2mm,bottom=2mm]
{\sffamily\bfseries\fontsize{22}{24}\selectfont\textcolor{#2}{#1}}\par
{\sffamily\bfseries\small #4}
\end{tcolorbox}
}

% ── Piezas del rediseno 2026 (legibilidad 6.º-11.º) ───────────────────
% Banda de estacion: reemplaza el titlebox macizo. Titulo a la izquierda,
% dato practico (tiempo/modalidad) a la derecha, en una sola linea.
% Nunca huerfana: pide 9 lineas libres (o salta de pagina rellenando con
% \vfill) y prohibe el salto justo despues de la banda. Nueve y no seis:
% la caja partible que sigue necesita ~80pt (titulo adosado, relleno y dos
% lineas) para arrancar en la misma pagina; con menos, tcolorbox la manda
% entera a la siguiente y la banda se queda sola. El penalty 10000 va
% en `after`, ANTES del espacio posterior: un `after skip` (aunque sea 0pt)
% es un \vskip tras la caja, y ese glue es un punto de corte legal que un
% \nopagebreak escrito despues de \end{tcolorbox} ya no cierra. Cada banda
% deja marcador PDF y actualiza la cabecera derecha.
\newcounter{milcestacion}
\newcommand{\milcestacion}[2]{%
\Needspace*{9\baselineskip}%
\stepcounter{milcestacion}%
\pdfbookmark[1]{#2}{milcestacion\themilcestacion}%
\markright{#2}%
\begin{tcolorbox}[enhanced,colback=#1,colframe=#1,arc=2mm,boxrule=0pt,
  left=5mm,right=5mm,top=2.6mm,bottom=2.6mm,before skip=11pt,
  after={\par\nopagebreak\vskip7pt}]
{\sffamily\bfseries\fontsize{15}{18}\selectfont\milctintasobre{#1}#2}
\end{tcolorbox}}

% Tarjeta de paso numerada: sustituye las tablas «Pilar / Aplicacion», que
% eran formato de consulta docente, no de estudiante. El numero va en un
% circulo sobre el borde para que la secuencia se lea de un vistazo.
\newcommand{\milcpaso}[3]{%
\Needspace{4\baselineskip}%
\begin{tcolorbox}[enhanced,colback=white,colframe=milcLinea,arc=1.5mm,boxrule=.9pt,
  left=14mm,right=4mm,top=3mm,bottom=3mm,before skip=4pt,after skip=4pt,
  fontupper=\RaggedRight,
  overlay={\node[anchor=north west,circle,fill=#2,text=white,inner sep=0pt,
    minimum size=7.4mm,font=\sffamily\bfseries\small]
    at ([xshift=3.6mm,yshift=-3.2mm]frame.north west) {#1};}]
#3
\end{tcolorbox}}

% Casilla marcable de tamano util con lapiz (la anterior era \fbox{\phantom{x}},
% de alto variable segun el texto que la seguia).
\newcommand{\milcmarca}{\framebox[3.8mm]{\rule{0pt}{3.4mm}}}

% Fila marcable de lista suelta (checks de la estacion 1).
\newcommand{\milccheckrow}[1]{%
\par\vspace{1.8mm}\noindent
\makebox[6.4mm][l]{\raisebox{-.55mm}{\textcolor{milcNegro}{\milcmarca}}}%
\parbox[t]{\dimexpr\linewidth-6.4mm\relax}{\RaggedRight #1}\par}

% Celda marcable dentro de la rubrica (tabla de autoevaluacion). Misma
% sangria francesa, pero pensada para vivir dentro de una columna X.
\newcommand{\milcopcion}[1]{%
\makebox[5.2mm][l]{\raisebox{-.55mm}{\textcolor{milcNegro}{\milcmarca}}}%
\parbox[t]{\dimexpr\linewidth-5.2mm\relax}{\RaggedRight #1}}

% ── Recursos visuales de la guía ──────────────────────────────────────
% Declarados en el bloque `recursos:` del YAML y emitidos por el builder.
% \guiaFigura   → imagen rasterizada o PDF (foto, carta celeste, montaje).
% \guiaDiagrama → fuente TikZ/circuitikz (.tex) incrustada como vector:
%                 es la vía para circuitos y esquemáticos editables.
% [#1] = texto alternativo (alt) · #2 = ruta relativa al .tex · #3 = pie
% El alt llega al PDF etiquetado (/Alt) y lo lee el lector de pantalla.
\newcommand{\guiaFigura}[3][]{%
\begin{center}
\includegraphics[width=0.88\linewidth,height=0.38\textheight,keepaspectratio,alt={#1}]{#2}\\[1.5mm]
{\footnotesize\itshape\textcolor{milcNegro!75}{#3}}
\end{center}
}

\newcommand{\guiaDiagrama}[2]{%
\begin{center}
\input{#1}\\[1.5mm]
{\footnotesize\itshape\textcolor{milcNegro!75}{#2}}
\end{center}
}

\begin{document}
\thispagestyle{empty}

% ── Portada ───────────────────────────────────────────────────────────
% Antes: un rectangulo de color a pagina completa. Se veia bien en pantalla,
% pero en la fotocopiadora del colegio se comia el toner y salia gris sucio.
% Ahora la banda ocupa el tercio superior y el resto va en blanco.
\begin{tikzpicture}[remember picture,overlay]
  \path[fill=milcGradePrimary]
    (current page.north west) rectangle ([yshift=-10.6cm]current page.north east);

  \node[anchor=north west,text=milcGradeText] at ([xshift=2.0cm,yshift=-1.55cm]current page.north west)
    {\sffamily\bfseries\fontsize{12}{14}\selectfont GRADO};
  \node[anchor=north west,text=milcGradeText,opacity=.85] at ([xshift=2.0cm,yshift=-2.35cm]current page.north west)
    {\sffamily\bfseries\fontsize{8.5}{10}\selectfont SERIE GUÍAS · TECNOLOGÍA E INFORMÁTICA};
  \node[anchor=north east,text=milcGradeText,opacity=.85,align=right] at ([xshift=-2.0cm,yshift=-1.55cm]current page.north east)
    {\sffamily\bfseries\fontsize{9}{12}\selectfont I.E. Sor María Juliana\\Cartago, Valle del Cauca};

  \node[anchor=west,text=milcGradeText] (gradonum) at ([xshift=1.85cm,yshift=-7.1cm]current page.north west)
    {\sffamily\bfseries\fontsize{112}{112}\selectfont 8};
  % El circulito de grado (°) solo aplica a las guias de grado («11°»); en
  % Bebras y Semillero el numero grande es un momento/modulo, no un grado.
  % Se posiciona relativo al nodo `gradonum`, no en coordenadas absolutas.
  \draw[milcGradeText,line width=1.6pt]
    ([xshift=2.0mm,yshift=-4.5mm]gradonum.north east) circle (.15cm);

  \node[anchor=west,text=milcGradeText,text width=11.4cm,align=left] at ([xshift=1.5cm,yshift=0.5cm]gradonum.east)
    {
     {\sffamily\bfseries\fontsize{9}{11}\selectfont\MakeUppercase{Período 2 · Lógica, algoritmos y computación física}}\\[3mm]
     {\sffamily\bfseries\fontsize{21}{25}\selectfont\setlength{\baselineskip}{27pt}Decir qué funciona\\[4pt]y qué no}\\[3.5mm]
     {\sffamily\bfseries\fontsize{15}{18}\selectfont Guía 20-8-TIC}
    };
\end{tikzpicture}

\vspace*{9.4cm}
\vfill

{\sffamily\bfseries\fontsize{8}{10}\selectfont\textcolor{milcGradeDark}{LO QUE TE LLEVAS HOY}}

\begin{tcolorbox}[enhanced,colback=white,colframe=milcNegro,arc=2mm,boxrule=1.6pt,
  left=5mm,right=5mm,top=4mm,bottom=4mm,before skip=3pt,after skip=12pt,
  fontupper=\RaggedRight\fontsize{11}{15.5}\selectfont]
Una sustentación de cinco minutos del proyecto de monitoreo de la sesión 9, con cinco láminas, el micro:bit funcionando en vivo, dos cosas que funcionaron y una que todavía falla con su explicación, y de tres a cinco preguntas del grupo respondidas. Y una autoevaluación escrita con una fortaleza y una mejora concretas.
\end{tcolorbox}

\begin{tcolorbox}[enhanced,colback=milcGris,colframe=milcGris,arc=2mm,boxrule=0pt,
  left=5mm,right=5mm,top=3.5mm,bottom=3.5mm,after skip=12pt]
{\sffamily\bfseries\fontsize{8}{10}\selectfont\textcolor{milcNegro!55}{ESTA GUÍA ES DE}}\\[3mm]
\begin{tabularx}{\linewidth}{X p{3.2cm} p{3.2cm}}
\linead{\linewidth} & \linead{3.0cm} & \linead{3.0cm}\\[-1mm]
{\fontsize{8.5}{10}\selectfont\textcolor{milcNegro!55}{Nombre completo}} &
{\fontsize{8.5}{10}\selectfont\textcolor{milcNegro!55}{Grupo}} &
{\fontsize{8.5}{10}\selectfont\textcolor{milcNegro!55}{Fecha}}\\
\end{tabularx}
\end{tcolorbox}

\vfill

\noindent\textcolor{milcLinea}{\rule{\linewidth}{1pt}}\\[2mm]
{\sffamily\bfseries\fontsize{9.5}{12}\selectfont Autor: Dr. Álvaro Cárdenas Orozco}\\[1mm]
{\sffamily\fontsize{8.5}{11}\selectfont\textcolor{milcNegro!55}{Metodología MILC · apertura ancestral · pensamiento computacional · triángulo Dussel–estoicismo–Floridi}}

\newpage

\section*{Apertura: Sustentación del proyecto --- decir qué funciona y qué no}

\begin{softbox}{milcGradeDark}{milcGradeSoft}{Saber ancestral}
En la sesión 3 del periodo 1 viste que los pueblos indígenas del Cauca no escriben sus normas desde una oficina. Salen de un congreso: la gente de los resguardos llega, discute durante días y de ahí nacen los \textbf{mandatos}. El XV Congreso del CRIC se reunió en Río Blanco, Sotará, del 25 al 30 de junio de 2017 (Consejo Regional Indígena del Cauca, 2017). Lo que importa hoy es lo que pasa después. La organización tiene que volver ante la misma gente y decir, en público, cuáles mandatos se cumplieron y cuáles no. Y la propia organización reporta que el cumplimiento es incompleto. No lo esconde: lo dice. Un mandato no vale por estar escrito; vale porque quienes lo hicieron vuelven a mirarlo. La \textbf{cara de exclusión}: un mandato es un acto de gobierno de un pueblo que se disputa con el Estado y con actores armados; no es una dinámica de convivencia escolar. Lo que tú haces hoy es mucho más pequeño y aprende de eso: presentar tu proyecto ante tu grupo y decir, con los datos en la mano, qué funcionó y qué todavía no.\par\smallskip{\footnotesize\textcolor{milcNegro!70}{\textbf{Fuente.} Consejo Regional Indígena del Cauca. (2017). \emph{Avanza el XV congreso regional del CRIC en el resguardo de Rioblanco Sotará del pueblo indígena yanacona}.}}
\end{softbox}

\begin{softbox}{milcTurquesa}{milcGris}{Saber contemporáneo}
Una \textbf{sustentación técnica} es la presentación de un proyecto ante personas que pueden preguntar. En cinco minutos muestras el problema, la solución, los datos y la propuesta. Tiene tres condiciones. (1) \textbf{Demo en vivo}: el micro:bit funcionando frente al grupo, no una captura de pantalla. (2) \textbf{Limitaciones declaradas}: una cosa que aún falla, con su explicación técnica. (3) \textbf{Respuestas honestas}: si alguien encuentra un problema, se reconoce; si no sabes, dices «no lo resolví» y por qué. En ingeniería ningún sistema está libre de errores. La diferencia está entre quien los declara y quien los esconde. Un equipo que dice «el umbral de luz dispara falsas alarmas cuando pasa una nube» es más confiable que uno que dice «todo funciona». El guion de cinco minutos son unas 700 palabras. Se ensaya con cronómetro al menos dos veces.
\end{softbox}

\begin{softbox}{milcMostaza}{milcGris}{Pregunta-puente}
\textit{El CRIC vuelve ante su gente y dice cuáles mandatos no se cumplieron. Cuando muestres tu proyecto, ¿qué te cuesta más: enseñar lo que funciona o decir en voz alta lo que todavía falla?}
\end{softbox}

\begin{softbox}{milcVerde}{milcGris}{Saber y saber hacer}
Al terminar podrás: (1) \textbf{analizar} tu proyecto con honestidad y separar lo que funciona de lo que falla; (2) \textbf{crear} cinco láminas y un guion de cinco minutos con demo en vivo; (3) \textbf{evaluar} tu sustentación y la de otros con criterios observables; (4) \textbf{explicar} por qué declarar una limitación hace más confiable un proyecto.
\end{softbox}



\small
\begin{tabularx}{\textwidth}{>{\bfseries}p{3.0cm}Y}
\rowcolor{milcGradePrimary!12}Autor & Dr. Álvaro Cárdenas Orozco\\
DBA & Pensamiento computacional aplicado a sistemas de control con sensores y actuadores (MEN, T\&I)\\
Referentes & Pensamiento computacional · MakeCode / micro:bit · Logos como razón universal\\
Producto final & Una sustentación de cinco minutos del proyecto de monitoreo de la sesión 9, con cinco láminas, el micro:bit funcionando en vivo, dos cosas que funcionaron y una que todavía falla con su explicación, y de tres a cinco preguntas del grupo respondidas. Y una autoevaluación escrita con una fortaleza y una mejora concretas.\\
\end{tabularx}
\normalsize

\puente{En la sesión 9 dejaste el micro:bit midiendo tres jornadas y encontraste un patrón. Hoy lo presentas ante el grupo con el micro:bit en vivo y dices qué funciona y qué no. Con esto cierras el periodo 2; en el periodo 3 empieza el diseño de piezas multimedia.}

\milcestacion{milcGradeDark}{Ruta MILC: así vamos a aprender}
\begin{tabularx}{\textwidth}{>{\centering\arraybackslash}X>{\centering\arraybackslash}X>{\centering\arraybackslash}X>{\centering\arraybackslash}X}
\phasecard{1}{milcVerde}{milcGris}{Escucha\\[-1mm]\small Observo, escucho y pregunto.} &
\phasecard{2}{milcTurquesa}{milcGris}{Entender\\[-1mm]\small Sistematizo y ordeno.} &
\phasecard{3}{milcMagenta}{milcGris}{Hacer\\[-1mm]\small Construyo y aplico.} &
\phasecard{4}{milcVino}{milcGris}{Cierre\\[-1mm]\small Evalúo y reflexiono.}
\end{tabularx}


\puente{Antes de armar láminas, vas a hacer el inventario honesto de tu proyecto: tres cosas que funcionaron, dos que fallan. Es el reporte de cumplimiento, en pequeño.}

\milcestacion{milcVerde}{Estación 1 de 3 · Escucha}

\begin{softbox}{milcVerde}{milcGris}{Escena para escuchar}
\textbf{Actividad 1 · ANALIZA --- El inventario honesto} (15 min · individual). (1) Abre tu bitácora de treinta filas y tu programa de la sesión 9. (2) Escribe tres cosas que funcionaron: la lectura del sensor, la calibración, la alerta visible, la frecuencia cumplida, la bitácora completa. (3) Escribe dos cosas que fallan o son frágiles: el umbral que dispara falsas alarmas, las horas que no se cumplieron, las filas perdidas, la batería. (4) Para cada falla, escribe en una línea por qué crees que pasa. (5) Elige la falla que vas a declarar en la sustentación y escribe cómo la mejorarías.
\end{softbox}

{\sffamily\bfseries\fontsize{8}{10}\selectfont\textcolor{milcNegro!55}{MARCA CUANDO LO HAYAS HECHO}}
\milccheckrow{Tengo tres cosas que funcionaron, con evidencia en la bitácora.}
\milccheckrow{Tengo dos fallas con su porqué en una línea.}
\milccheckrow{Elegí la falla que voy a declarar y cómo la mejoraría.}

\begin{infoband}{milcVerde}


\textbf{Lo que va al cuaderno (Actividad 1 · ANALIZA).} \textbf{Título:} «Actividad 1 --- El inventario honesto». \textbf{Formato:} dos columnas: «funciona» con tres filas y «falla» con dos, cada una con una línea de porqué; y la falla elegida con su mejora. \textbf{Extensión:} media página. \textbf{Sabes que terminaste cuando} cada cosa que funciona se puede señalar en la bitácora, y cada falla tiene una causa y no una excusa.
\end{infoband}

\puente{Ya tienes el inventario. Ahora aprendes qué va en cada lámina, armas las cinco con tu pareja y ensayas con cronómetro.}

\milcestacion{milcTurquesa}{Estación 2 de 3 · Entender}

\begin{softbox}{milcTurquesa}{milcGris}{Lo que vas a entender}
\textbf{Actividad 2 · CREA --- Cinco láminas y un ensayo} (30 min · parejas). (1) Reduce tu proyecto a cinco láminas de máximo veinticinco palabras cada una: problema y variable; sensor, umbrales y alerta; los datos y el patrón con su cifra; la propuesta de acción; lo que falla y cómo se mejoraría. (2) Escribe el guion hablado: unas 700 palabras son cinco minutos. (3) Prepara la demo: micro:bit con batería, código cargado, y una forma de producir la alerta en vivo, tapar el sensor o calentarlo con la mano. (4) Ensaya frente a tu pareja con cronómetro. (5) Tu pareja te hace dos preguntas difíciles y te dice qué no entendió; cambien de rol.
\end{softbox}

\milcpaso{1}{milcTurquesa}{\textbf{Las cinco láminas y su tiempo.} Problema y variable, 45 segundos. Sensor, umbrales y alerta, con una foto del micro:bit instalado, 60 segundos. Datos y patrón, con un gráfico simple de las treinta filas, 75 segundos. Propuesta, 45 segundos. Lo que falla y cómo se mejora, 45 segundos. Quedan unos segundos para respirar.}
\milcpaso{2}{milcTurquesa}{\textbf{Las preguntas que van a llegar.} «¿Cómo elegiste los umbrales?»: con la bitácora de calibración. «¿Y si el sensor se daña?»: el sistema queda mudo, y lo dices. «¿Por qué solo treinta mediciones?»: defiendes la muestra o reconoces el límite. «¿Qué harías con otra semana?»: respondes con la mejora que ya escribiste.}
\milcpaso{3}{milcTurquesa}{\textbf{La demo en vivo.} La captura de pantalla no demuestra nada; el micro:bit mostrando la alerta cuando tapas el sensor sí. Si algo impide la demo, sin batería, sin computador, lo dices al empezar y muestras el video que grabaste como respaldo.}
\milcpaso{4}{milcTurquesa}{\textbf{Cómo se hace, paso a paso.} (1) Cinco láminas cortas. (2) Guion de 700 palabras. (3) Demo preparada y probada. (4) Dos ensayos con cronómetro. (5) Preguntas de tu pareja. (6) Ajustar lo que no se entendió.}

\begin{drawbox}{milcTurquesa}{La sustentación, en una mirada}
\textbf{Problema:} qué mediste y por qué. \\[1mm] \textbf{Solución:} sensor, umbrales, alerta. \\[1mm] \textbf{Datos:} el patrón y su cifra. \\[1mm] \textbf{Propuesta:} qué se haría con el dato. \\[1mm] \textbf{Límite:} lo que falla y cómo se mejora.
\end{drawbox}

\begin{softbox}{milcMostaza}{milcGris}{Errores comunes y su corrección}
(1) \textbf{Leer las láminas}: el grupo lee más rápido que tú. (2) \textbf{Esconder la falla}: alguien la encuentra en las preguntas y pierdes más. (3) \textbf{Inflar el patrón}: «siempre» donde los datos dicen «dos de tres días». (4) \textbf{Captura en vez de demo}: nadie vio el micro:bit funcionar. (5) \textbf{Evadir la pregunta dura}: responder otra cosa en vez de «no lo resolví».
\end{softbox}

\begin{infoband}{milcTurquesa}


\textbf{Lo que va al cuaderno (Actividad 2 · CREA).} \textbf{Título:} «Actividad 2 --- Cinco láminas y un ensayo». \textbf{Formato:} el texto de las cinco láminas, el guion de 700 palabras, el tiempo de los dos ensayos y las dos preguntas de tu pareja con lo que no entendió. \textbf{Extensión:} una página. \textbf{Sabes que terminaste cuando} el ensayo cabe en cinco minutos y la lámina de lo que falla tiene una causa y una mejora.
\end{infoband}

\puente{Ya ensayaste. Toca sustentar frente al grupo con el micro:bit funcionando, responder preguntas y escribir tu autoevaluación.}

\milcestacion{milcMagenta}{Estación 3 de 3 · Hacer}

\begin{softbox}{milcMagenta}{milcGris}{Cómo vas a construir tu producto}
\textbf{Actividad 3 · EVALÚA --- Sustentar con el micro:bit en vivo} (35 min · individual). (1) Sustenta frente al grupo, o en rondas de cinco, con cronómetro visible: cinco minutos. (2) Haz la demo: produce la alerta en vivo tapando o calentando el sensor. (3) Declara la falla que elegiste antes de que te la pregunten. (4) Responde de tres a cinco preguntas; si no sabes, di «no lo resolví» y por qué. (5) Mientras escuchas a otros, evalúa con los seis criterios de la tabla y escribe una pregunta para cada uno. (6) Al final, escribe tu autoevaluación: una fortaleza y una mejora concretas.
\end{softbox}

\milcpaso{1}{milcMagenta}{\textbf{Tu entrega tiene cuatro piezas.} (1) Las cinco láminas. (2) El guion. (3) La lista de lo que funciona y lo que falla. (4) La autoevaluación con una fortaleza y una mejora. El programa y la bitácora ya los entregaste en la sesión 9.}
\milcpaso{2}{milcMagenta}{\textbf{Declarar antes de que pregunten.} Si dices «el umbral de luz falla cuando pasa una nube» en la lámina cinco, el grupo lo recibe como honestidad. Si lo callas y alguien lo descubre en las preguntas, el grupo lo recibe como algo que escondiste.}
\milcpaso{3}{milcMagenta}{\textbf{Responder con datos.} «¿Tu propuesta es factible?» se responde con la cifra: «subió de 24 a 28 los tres días, abrir ventanas a las 11:30 cuesta cero». Sin cifra, es una opinión contra otra.}
\milcpaso{4}{milcMagenta}{\textbf{La autoevaluación concreta.} Fortaleza: «la demo salió a la primera porque la probé dos veces antes». Mejora: «me pasé 40 segundos en la lámina de datos; la próxima recorto el gráfico». No «hablar mejor» ni «prepararme más».}

\begin{drawbox}{milcMagenta}{Antes de entregar}
\checkbox Cinco láminas de máximo veinticinco palabras cada una. \\[1mm] \checkbox Guion de unas 700 palabras ensayado dos veces con cronómetro. \\[1mm] \checkbox Demo en vivo hecha, o el motivo dicho al empezar y el video de respaldo. \\[1mm] \checkbox La falla declarada antes de las preguntas, con causa y mejora. \\[1mm] \checkbox De tres a cinco preguntas respondidas, con «no lo resolví» cuando tocó. \\[1mm] \checkbox Sustentación dentro de los cinco minutos. \\[1mm] \checkbox Autoevaluación con una fortaleza y una mejora concretas.
\end{drawbox}

\begin{softbox}{milcMagenta}{milcGris}{Plantilla de trabajo}
\textbf{Problema.} \textit{Medí … porque …} \\[1mm] \textbf{Solución.} \textit{Con el sensor de …, umbrales de … y …, y alerta …} \\[1mm] \textbf{Datos.} \textit{Los tres días, … entre … y …: de … a …} \\[1mm] \textbf{Propuesta.} \textit{Por eso propongo …} \\[1mm] \textbf{Límite.} \textit{Todavía falla … porque …; lo mejoraría con …}
\end{softbox}

\begin{infoband}{milcMagenta}


\textbf{Lo que va al cuaderno (Actividad 3 · EVALÚA).} \textbf{Título:} «Actividad 3 --- Sustentar con el micro:bit en vivo». \textbf{Formato:} las preguntas que te hicieron con tu respuesta en una línea, tu evaluación de dos compañeros con los seis criterios, y tu autoevaluación. \textbf{Extensión:} una página. \textbf{Sabes que terminaste cuando} la fortaleza y la mejora son cosas que alguien más pudo ver en tu sustentación.
\end{infoband}

\puente{Lo que entregas hoy: las cinco láminas, el guion, la lista de lo que funciona y lo que falla, y la autoevaluación. La prueba de calidad: el micro:bit funcionó en vivo y dijiste una limitación sin que te la preguntaran.}

\milcestacion{milcMostaza}{Producto de la sesión}
\textbf{Sustentación de cinco minutos + cinco láminas + demo en vivo + falla declarada + preguntas respondidas + autoevaluación}.
\begin{softbox}{milcMostaza}{milcGris}{Criterios de calidad}
(1) Las cinco láminas tienen problema, solución, datos con cifra, propuesta y límite. (2) El micro:bit funcionó en vivo, o se dijo el motivo al empezar y hubo video de respaldo. (3) La falla se declaró antes de las preguntas, con causa y mejora. (4) Las preguntas se respondieron con datos o con «no lo resolví» y su porqué. (5) La sustentación cupo en cinco minutos. (6) La autoevaluación tiene una fortaleza y una mejora que otros pudieron ver.
\end{softbox}

\puente{Antes de cerrar, mira lo que hiciste desde cinco lados. Una sustentación que declara una falla enseña más que una que la esconde y la audiencia la descubre.}

\milcestacion{milcVino}{Evaluación liberadora · cinco dimensiones}

\begin{softbox}{milcVino}{milcGris}{Sentido de la evaluación}
Evaluar no es solo revisar una nota. Es reconocer qué aprendí, qué cambió en mí, qué emociones manejé, cómo me relaciono con la ciudad y la comunidad, y qué le dejaré a quien viene detrás.
\end{softbox}

\textbf{Mi reflexión liberadora en cinco dimensiones.} Completa cada frase iniciada:

\small
\begin{tabularx}{\textwidth}{>{\bfseries\RaggedRight\arraybackslash}p{3.4cm}Y>{\RaggedRight\arraybackslash}p{4.6cm}}
\rowcolor{milcVino}\color{white}Dimensión & \color{white}\bfseries Mi respuesta (frame starter) & \color{white}\bfseries Algo que descubrí\\
1. Personal & Lo que más cambió en mí fue \linead{6.5cm} & \linead{4.2cm}\\
2. Emocional & La emoción más fuerte fue \linead{2.5cm} y la regulé \linead{3cm} & \linead{4.2cm}\\
3. Ciudadana & En democracia esto importa porque \linead{6cm} & \linead{4.2cm}\\
4. Local (Cartago) & En mi barrio sería útil para \linead{6.5cm} & \linead{4.2cm}\\
5. Intergeneracional & A mi abuela/abuelo le diría \linead{6.5cm} & \linead{4.2cm}\\
\end{tabularx}
\normalsize

\begin{infoband}{milcVino}
\textbf{Para la próxima sustentación.} Vuelve a esta tabla en seis meses. Verás que las respuestas cambian — no porque lo aprendido cambie, sino porque tú cambias. La evaluación liberadora es una conversación contigo a través del tiempo.
\end{infoband}


\puente{Para terminar, tres ideas sobre presentarse ante otros, contadas con palabras de esta guía: Dussel sobre el rostro que interpela, Marco Aurelio sobre ser bueno en vez de discutirlo, y Floridi sobre la mitad de los datos que no sirve.}

\milcestacion{milcGradeDark}{Tres ideas para cerrar · Dussel, un estoico y Floridi}

\begin{softbox}{milcVino}{milcGris}{Enrique Dussel · lente del nosotros}
\textit{El otro se muestra como otro cuando, desde fuera de mi sistema, me interpela con su rostro y su reclamo.}

{\footnotesize\itshape\color{milcNegro!70}--- idea tomada de Enrique Dussel, contada con palabras de esta guía. No es una cita textual.}

Tu proyecto es tu sistema. Las preguntas del grupo vienen de fuera y te muestran lo que desde adentro no viste. Responder con honestidad es dejar que el otro te interpele, no defenderte.

\textbf{¿Qué pregunta de hoy me mostró algo de mi proyecto que yo no había visto?}
\end{softbox}
\linead{\linewidth}\\[1mm]
\linead{\linewidth}

\begin{softbox}{milcMostaza}{milcGris}{Estoicismo · Marco Aurelio · cuidado interior}
\textit{Deja de discutir cómo debería ser una buena persona y procura serlo.}

{\footnotesize\itshape\color{milcNegro!70}--- idea tomada de Marco Aurelio, contada con palabras de esta guía. No es una cita textual.}

Decir «mi sistema funciona» es discutir. Tapar el sensor y que salga la alerta frente al grupo es serlo. La demo en vivo vale más que cualquier lámina que lo afirme.

\textbf{¿Qué afirmé hoy con palabras que habría podido mostrar con el micro:bit?}
\end{softbox}
\linead{\linewidth}\\[1mm]
\linead{\linewidth}

\begin{softbox}{milcTurquesa}{milcGris}{Luciano Floridi · lente de la infoesfera}
\textit{La mitad de nuestros datos no sirve; el problema es que no sabemos cuál mitad.}

{\footnotesize\itshape\color{milcNegro!70}--- idea tomada de Luciano Floridi, contada con palabras de esta guía. No es una cita textual.}

De tus treinta filas, algunas tienen huecos, horas corridas o una nube que pasó. Declararlas es decir cuál mitad no sirve. Un proyecto que sabe cuáles datos son flojos es más confiable que uno que los presenta todos iguales.

\textbf{¿Cuáles de mis treinta filas no me atrevería a defender, y lo dije?}
\end{softbox}
\linead{\linewidth}\\[1mm]
\linead{\linewidth}

\begin{infoband}{milcNegro}
\textbf{Nota para el docente.} Las tres ideas están contadas con palabras de esta guía a partir del pensamiento de cada autor; no son citas textuales. De 9.º en adelante se trabajan con la obra y su referencia.
\end{infoband}

\milcestacion{milcVino}{Mi autoevaluación · marca una casilla por fila}

{\small
\setlength{\extrarowheight}{2.2pt}
\begin{tabularx}{\textwidth}{>{\RaggedRight\arraybackslash\bfseries}p{3.0cm}YYY}
\rowcolor{milcVino}\color{white}\bfseries Criterio & \color{white}\bfseries Lo logré & \color{white}\bfseries En proceso & \color{white}\bfseries Necesito apoyo\\[.6mm]
Comprensión del tema & \milcopcion{Explico las ideas con mis palabras.} & \milcopcion{Reconozco las ideas, falta precisión.} & \milcopcion{Necesito repasar lo básico.}\\[.8mm]
\rowcolor{milcGris}Pensamiento computacional & \milcopcion{Usé los pilares al armar mi producto.} & \milcopcion{Usé algunos pilares.} & \milcopcion{Necesito releer la estación 2.}\\[.8mm]
Aplicación y producto & \milcopcion{Mi producto cumple los criterios.} & \milcopcion{Está armado, con ajustes pendientes.} & \milcopcion{Aún está en construcción.}\\[.8mm]
\rowcolor{milcGris}Reflexión 5 dimensiones & \milcopcion{Respondí cada una con honestidad.} & \milcopcion{Respondí algunas con detalle.} & \milcopcion{Mis respuestas son muy generales.}\\[.8mm]
Triángulo de pensamiento & \milcopcion{Conecté las 3 voces con mi trabajo.} & \milcopcion{Conecté al menos una.} & \milcopcion{No las relaciono con mi proceso.}\\[.8mm]
\end{tabularx}}

\vspace{3mm}
\textbf{Hoy entendí que:}\\[1mm]
\linead{\linewidth}\\[1mm]
\linead{\linewidth}

\textbf{Mi compromiso para la próxima sesión es:}\\[1mm]
\linead{\linewidth}\\[1mm]
\linead{\linewidth}

\begin{softbox}{milcMostaza}{milcGris}{Cita de despedida · Marco Aurelio}
\textit{``El obstáculo en el camino se vuelve el camino. Lo que estorba al camino se vuelve camino.''}\\[1mm]
Cada error de hoy, bien observado, se convierte en pieza del camino que pisarás mañana.
\end{softbox}

\small
\begin{tabularx}{\textwidth}{>{\bfseries}p{3.6cm}Y}
\rowcolor{milcGradeDark!12}Nombre completo: & \linead{10cm}\\
Fecha: & \linead{4cm} \quad Curso/grupo: \linead{4cm}\\
Mi firma: & \linead{9cm}\\
\end{tabularx}
\normalsize

\begin{infoband}{milcGradeDark}
\textbf{Para el periodo 3.} Con esta sustentación cierras «Lógica y micro:bit». Guarda el programa, la bitácora y las láminas: son tu evidencia del periodo. En la sesión 1 del periodo 3 vas a empezar con diseño visual y jerarquía de información en piezas multimedia. Las cinco láminas de hoy son tu primer borrador de eso.
\end{infoband}

\end{document}
